\documentclass[UTF8,a4paper,zihao=-4,fontset=fandol]{ctexart}
% 版心四边留白 2.85 cm，保证实际着墨边界 ≥2.5 cm（规范第一条）
\usepackage[left=2.85cm,top=2.85cm,bottom=2.85cm,right=2.85cm,includefoot]{geometry}
\usepackage{booktabs,array,longtable,enumitem,fancyhdr}
\setlength{\parindent}{2em}
\setlist{leftmargin=2em,itemsep=2pt,topsep=3pt}
\linespread{1.04}\selectfont
\pagestyle{fancy}\fancyhf{}\fancyfoot[C]{\thepage}\renewcommand{\headrulewidth}{0pt}
\begin{document}
\begin{center}{\zihao{3}\bfseries AI工具使用详情}\end{center}

\section*{一、工具名称、版本或型号}
竞赛过程中使用了以下三款大语言模型工具：Kimi K3、GLM 5.3 及 DeepSeek V4.1 Flash。本文所称AI工具仅指上述交互式问答与辅助工具。

需要明确说明：本题的\textbf{主体建模工作与全部编程实现均由参赛队员人工完成}，包括集合几何建模、七点发现环与21点双环布点、剩余覆盖不变量与固定点数下界的推导、控制器与在线程序的编写调试、实验方案设计、以及模拟器正式测试的启动与确认。AI工具在竞赛过程中仅起辅助作用，未主导上述任何核心环节；数学建模方案的取舍、正式测试的启动与确认、最终成绩核对及材料提交也均由参赛队员完成。

\section*{二、具体使用目的和环节}
\begin{longtable}{>{\raggedright\arraybackslash}p{3.0cm}>{\raggedright\arraybackslash}p{9.9cm}}
\toprule 使用环节&具体用途\\\midrule\endhead
代码调试&协助排查HTTP请求字段、串行请求与幂等请求号的实现问题；定位编译与运行异常；整理可复现运行入口与依赖说明。所有修改建议均由队员在本地编译并运行验证后才采纳。\\
数值检验&按队员给定的复算口径，独立核对成功清除次数、虚拟总时间、平均定位清除时间、程序运行时间与HTTP错误数；交叉核对本地配对实验、源数--单源均时附加分析与论文数表之间的一致性。\\
论文语言严谨性润色&压缩重复表述，统一术语与符号，区分“几何可行性证书”“本地统计结果”和“官方正式测试记录”的证据边界，删除无法由题面、代码、实验表或官方日志支持的强结论，避免把未知正式案例的主动退出误写为已知真值下的100\%清除。\\
\bottomrule
\end{longtable}

\section*{三、主要提示方式与使用过程}
采用“队员给出任务与本地文件、要求先核验再修改”的提示方式。典型交互经过概括后如下：
\begin{enumerate}
\item “读取当前B题题面、论文、题解和模拟器程序，先在演练测试中验证，不触碰次数受限的正式测试。”
\item “根据多次问题三、问题四演练和正式测试日志，汇总清除数、平均定位清除时间和程序运行时间，并与本地实验比较。”
\item “依据2026年论文格式规范和AI工具使用规定修改论文，补齐正式测试表、模型优缺点与推广、AI工具使用声明，并整理匿名支撑材料。”
\item “检查支撑包是否含身份信息，保留官方导出的加密日志文件名，生成文件清单和校验值。”
\end{enumerate}
AI工具先读取题面、协议、官方规范和已有程序产物，再提出修改建议或生成修改稿；涉及正式测试时，由队员在模拟器中人工开启窗口并明确确认后才运行程序。AI工具没有代替队员获取隐藏案例真值，也没有修改官方加密日志。

\section*{四、采纳、人工修改和核验情况}
\begin{enumerate}
\item \textbf{模型与算法：}核心集合几何、七点发现环、21点双环、剩余覆盖证书、有限动作规划和光学兜底均由参赛队员建立与推导。AI工具主要负责检查论证衔接、实现一致性和结果表达；是否保留候选策略由本地配对实验和正式测试风险共同决定。
\item \textbf{程序：}在线控制器与实验程序的编写、调试和参数配置由队员完成。AI工具提出的代码检查建议均先通过本地回归测试、公开协议模拟器或不限次数的演练测试，再由队员决定是否用于正式测试。正式测试使用冻结配置，六次均由模拟器生成加密日志并进入上传队列。
\item \textbf{数值：}论文中的六次正式成绩由程序日志逐条复算，并与模拟器导出的案例编码核对。平均定位清除时间按“虚拟总时间/成功清除数”计算；因正式测试不公开干扰源总数，论文没有虚构清除比例。所有复算结果均由队员对照原始日志二次确认。
\item \textbf{论文：}队员逐段审阅AI工具整理的文字，核对公式、参数、表格、案例编码和结论；删除无法由题面、代码、实验表或官方日志支持的强结论，并补充局限性和适用边界。
\item \textbf{提交材料：}队员最终检查论文与支撑材料的匿名性、文件大小、日志原文件名、PDF可读性和压缩包目录结构后再提交。
\end{enumerate}

\section*{五、责任说明}
参赛队对所提交论文、程序、数据处理、实验结果和结论承担全部责任。AI工具输出仅作为辅助材料，未经人工审查、运行验证或证据核对的内容未直接作为核心建模与分析成果提交。
\end{document}
